% --- Άσκηση 33698 (33698.tex) ---
\Askhsh{33698}{4}
\begin{ylh}{2}
    \item 3.3. Εξισώσεις 2ου Βαθμού
    \item 4.2. Ανισώσεις 2ου Βαθμού
\end{ylh}
Δίνεται το τριώνυμο $f(x) = x^{2} - 6x + \lambda - 3$, με $\lambda\mathbb{\in R\ .}$

\begin{alist}
\item Να υπολογίσετε την διακρίνουσα $\Delta$ του τριωνύμου.
\monades{5}
\item Να βρείτε τις τιμές του $\lambda$ για τις οποίες το τριώνυμο έχει δύο άνισες πραγματικές ρίζες.
\monades{7}
\item Αν $3 < \lambda < 12$ τότε:


\begin{rlist}
\item
  Να δείξετε ότι το τριώνυμο έχει δύο άνισες θετικές ρίζες.
\monades{6}
\item
  Αν $x_{1},x_{2}$ με $x_{1} < x_{2}$ είναι οι δύο ρίζες του τριωνύμου και $\kappa$, $\mu$ είναι δύο αριθμοί με $\kappa < 0$ και $x_{1} < \mu < x_{2}$, να προσδιορίσετε το πρόσημο του γινομένου\\
  $\kappa \cdot f(\kappa)\text{⋅}\mu \cdot f(\mu)$. Να αιτιολογήσετε την απάντησή σας
\monades{7}
\end{rlist}
\end{alist}
